\documentclass[11pt]{article}
\usepackage[margin=2.5cm]{geometry}
\usepackage{amsmath,amssymb,amsthm}
\usepackage{mathtools}
\usepackage{hyperref}

\newtheorem{theorem}{Theorem}
\newtheorem{lemma}[theorem]{Lemma}
\newtheorem{proposition}[theorem]{Proposition}
\newtheorem{corollary}[theorem]{Corollary}
\newtheorem{definition}[theorem]{Definition}
\newtheorem{remark}[theorem]{Remark}

\newcommand{\R}{\mathbb{R}}
\newcommand{\Q}{\mathbb{Q}}
\newcommand{\N}{\mathbb{N}}
\newcommand{\iv}[1]{[\![#1]\!]}
\DeclareMathOperator{\supp}{supp}

\title{From Probabilistic Population Protocols to \\
Non-Autocatalytic Reactions: An Informal Walkthrough}
\author{}
\date{2026-04-24}

\begin{document}
\maketitle

\begin{abstract}
This note walks through the proof of BD Theorem~11.2 (PP~$\to$~NAP) in
ordinary mathematical prose. No formal-verification language appears.
The goal is for the reader to be able to evaluate the mathematical
content on its own terms, step by step, before consulting any machine
artifact.
\end{abstract}

\section{What we want to prove}

A \emph{probabilistic population protocol} (PP) on $n$ species is
specified by a tensor of non-negative rationals
\[
   c_{r,i,j} \in \Q_{\ge 0}, \qquad r,i,j \in \{1,\dots,n\},
\]
symmetric in $(i,j)$ and subject to the mass-conservation constraint
\begin{equation}
   \sum_{r=1}^n c_{r,i,j} \;=\; 2 \qquad \text{for every pair } (i,j).
   \tag{P1}
\end{equation}
The reading is: when one copy of species $i$ collides with one copy of
species $j$, they react into a pair of output copies, with $c_{r,i,j}/2$
being the expected number of $r$-copies produced. Condition (P1) records
that two input copies give exactly two output copies (mass is conserved).

On the probability simplex $\Delta^{n-1} = \{x \in \R_{\ge 0}^n : \sum_k x_k = 1\}$,
the associated mean-field ODE is
\begin{equation}
   \dot x_r \;=\; \sum_{i,j} c_{r,i,j}\, x_i x_j \;-\; 2 x_r.
   \tag{P2}
\end{equation}
The quadratic term is the \emph{production} rate of $r$; the $-2x_r$ is the
\emph{degradation} rate, which integrates the fact that whenever any
reaction happens $r$ loses a copy proportional to its concentration.

A \emph{non-autocatalytic reaction} (NAP) in the chemical sense is an
elementary reaction $I \to O$ where neither $I$ nor $O$ dominates the
other by pure replication of the same species. Formally, we will work
degree-$3$: encode states as cubic multinomials $v_\alpha = x^\alpha$ with
$|\alpha| = 3$, and ask that the chain-rule rate of each $\dot v_\alpha$
factor into a sum of transitions $\beta \to \gamma$ of cubed exponents
where $\beta \ne \alpha$ and $\gamma \ne \alpha$ (``no species just
catalyzes itself''). The precise condition is the \emph{balanced
non-autocatalytic split} defined in \S\ref{sec:per-pair}.

\medskip

\textbf{The claim.} Every PP can be implemented by a collection of NAPs
in cubed coordinates, after adding one slack species.

\section{Step 1: folding degradation into a signed quadratic field}

Define the signed coefficients
\begin{equation}
   q_{r,i,j} \;:=\; c_{r,i,j} - \iv{i=r} - \iv{j=r},
\end{equation}
where $\iv{\cdot}$ denotes the Iverson bracket. Using $\sum_k x_k = 1$
twice to rewrite $2x_r = 2 x_r \sum_k x_k = \sum_{i,j} (\iv{i=r}+\iv{j=r}) x_i x_j$,
we get on the simplex
\begin{equation}
   \dot x_r \;=\; \sum_{i,j} q_{r,i,j}\, x_i x_j.
   \tag{P3}
\end{equation}
The signed tensor $q$ inherits the conservation law
\begin{equation}
   \sum_r q_{r,i,j} \;=\; \Big(\sum_r c_{r,i,j}\Big) - 2 \;=\; 0.
   \tag{P4}
\end{equation}

So far (P3)--(P4) is a lossless rewriting of the PP. Call any tensor
$q_{r,i,j}$ satisfying (P4) a \emph{quadratic conservative field} (QCF).
Every PP is a QCF by folding; not every QCF comes from a PP. The passage
from PP to QCF is a purely algebraic book-keeping step.

\medskip

\textbf{Reading of the signs.} A positive entry $q_{r,i,j} > 0$ means
that monomial $x_i x_j$ contributes a net \emph{production} of $r$.
A negative entry means the corresponding monomial is a net
\emph{consumer} of $r$ --- but notice that by construction these
negative entries come only from the folding terms
$-\iv{i=r} - \iv{j=r}$, which are not free parameters: they encode
``whenever $i$ reacts with $j$, the species $i$ and the species $j$
are each spent once''. The magnitude of a consumer entry is tied by
(P4) to the production entries in its own column.

\begin{lemma}[no positive self-self on the diagonal]
\label{lem:nss}
For any PP, the folded tensor satisfies
\[
   q_{r,r,r} \;\le\; 0 \qquad \text{for every } r.
\]
\end{lemma}
\begin{proof}
From $c_{s,r,r} \ge 0$ and $\sum_s c_{s,r,r} = 2$ one has
$c_{r,r,r} \le \sum_s c_{s,r,r} = 2$. Therefore
$q_{r,r,r} = c_{r,r,r} - 2 \le 0$.
\end{proof}

We will denote this property by (NSS). It is automatic for any PP.

\section{Step 2: the $\lambda$-lift --- carving out a slack species}
\label{sec:lambda}

Pick a species $j_0$ and a parameter $\lambda \in (0,1)$. Introduce one
new species, which we call $0$ (the \emph{slack}), and rename
\[
   u := \lambda x_{j_0}, \qquad
   r_0 := (1-\lambda) x_{j_0}, \qquad
   u + r_0 = x_{j_0}.
\]
For every other species $i \ne j_0$, keep $y_i := x_i$. The enlarged
state is $(r_0, y_1, \dots, y_{j_0-1}, u, y_{j_0+1}, \dots, y_n)$. On the
simplex, $\sum_i y_i + u + r_0 = 1$. The scalar $\alpha_i := \lambda$ if
$i = j_0$ else $1$ records the local rescaling.

\textbf{Goal of the lift.} We want the enlarged field to satisfy two
regularity properties:
\begin{itemize}
\item[(NSM)] \emph{No Self-Mix on slack:} $\widetilde q_{r, 0, k} = 0$
   and $\widetilde q_{r, k, 0} = 0$ for every $r$ and $k$. In words: the
   slack species $0$ never appears as a reactant factor.
\item[(NSS)] $\widetilde q_{r,r,r} \le 0$ for every $r$ in the enlarged
   index set.
\end{itemize}
Their conjunction we call \emph{slack-structured}. The slack species'
role is purely to absorb stoichiometric excess: it will be consumed or
produced but never participate \emph{quadratically} as a reactant with
itself or with anyone else.

\textbf{Lift formula.} Substitute $x_{j_0} = (u + r_0)/\lambda \cdot \lambda
= u + r_0$ and enforce the slack dynamics
$\dot r_0 = \dot x_{j_0} - \dot u = (1-\lambda) \dot x_{j_0}$.
The signed coefficient in the enlarged tensor is
\begin{equation}
\widetilde q_{r',\,i',\,k'} \;=\;
\begin{cases}
   0 & \text{if } i' = 0 \text{ or } k' = 0,\\[2pt]
   \dfrac{(1-\lambda)\, q_{j_0, I, K}}{\alpha_I \alpha_K}
      & \text{if } r' = 0,\ i' = I^+,\ k' = K^+,\\[6pt]
   \dfrac{\alpha_R\, q_{R, I, K}}{\alpha_I \alpha_K}
      & \text{if } r' = R^+,\ i' = I^+,\ k' = K^+,
\end{cases}
\tag{$\Lambda$}
\end{equation}
where $R^+, I^+, K^+$ indicate the image of the old species indices in
the enlarged set (slot $0$ is slack; old index $j$ goes to $j^+ = j+1$
abstractly). The rescaling $1/(\alpha_I \alpha_K)$ compensates for the
fact that an old monomial $x_I x_K$ becomes, after substitution,
$u \cdot x_K / \lambda = \alpha_I \alpha_K \cdot \ldots$; the upshot is
that the vector field $\dot{(\cdot)}$ is preserved.

\begin{lemma}[$\widetilde q$ is still conservative]
\label{lem:lift-cons}
For every $(i',k')$, $\sum_{r'} \widetilde q_{r', i', k'} = 0$.
\end{lemma}
\begin{proof}
For $i' = 0$ or $k' = 0$ every summand is zero.
For $i' = I^+, k' = K^+$, summing the $r' = 0$ row and the $r' = R^+$ rows:
\[
\alpha_I \alpha_K \sum_{r'} \widetilde q_{r', I^+, K^+}
  \;=\; (1-\lambda)\, q_{j_0, I, K} + \sum_R \alpha_R\, q_{R, I, K}.
\]
Use $\alpha_R = 1$ for $R \ne j_0$ and $\alpha_{j_0} = \lambda$ to break
off the $R = j_0$ term:
\[
  = (1-\lambda)\, q_{j_0, I, K} + \sum_{R} q_{R, I, K} + (\lambda - 1)\, q_{j_0, I, K}
  = \sum_R q_{R, I, K} = 0,
\]
using (P4) for the original $q$.
\end{proof}

\begin{lemma}[$\widetilde q$ is slack-structured]
\label{lem:lift-slack}
If the original $q$ satisfies (NSS), and $\lambda > 0$, then
$\widetilde q$ satisfies both (NSM) and (NSS).
\end{lemma}
\begin{proof}
(NSM) is immediate from the first line of ($\Lambda$).
For (NSS) on the slack row, $\widetilde q_{0,0,0} = 0 \le 0$.
For a non-slack row $r' = R^+$:
\[
\widetilde q_{R^+, R^+, R^+} \;=\; \frac{\alpha_R\, q_{R,R,R}}{\alpha_R^2}
   \;=\; \frac{q_{R,R,R}}{\alpha_R} \;\le\; 0,
\]
since $q_{R,R,R} \le 0$ by hypothesis and $\alpha_R > 0$.
\end{proof}

\begin{remark}[why we also need $\lambda < 1$]
The algebraic lemma uses only $\lambda > 0$. The hypothesis
$\lambda < 1$ is used downstream for the simplex interpretation: the
slack coordinate $r_0 = (1-\lambda) x_{j_0}$ must be non-negative, so
$1 - \lambda \ge 0$, and strict positivity $1 - \lambda > 0$ keeps the
slack non-trivial.
\end{remark}

At the end of this step we have, from any PP, a QCF on $n+1$ species
(slot $0$ is slack) that satisfies slack-structured = (NSM) + (NSS).
This is the field we actually split.

\section{Step 3: the $r^2$-trick and cubic coordinates}

\subsection*{3a. $r^2$-trick: adding a slack factor to every production}

Multiply every production monomial by $x_0^2$. Concretely, replace
\[
   \widetilde q_{r,a,b}\, x_a x_b
   \;\longrightarrow\;
   \widetilde q_{r,a,b}\, x_0^2 x_a x_b.
\]
On the simplex, multiplying by $x_0^2 \le 1$ is a pure time rescaling;
the qualitative dynamics are unaffected. The payoff is that every
production monomial now carries \emph{two} guaranteed slack factors.
Writing $\mu_{0,a,b}(k) := 2\,\iv{k=0} + \iv{k=a} + \iv{k=b}$ for the
exponent vector of $x_0^2 x_a x_b$, this is the monomial that will be
split.

\subsection*{3b. Cubing: moving to degree-$3$ coordinates}

For an exponent vector $\alpha : \{0,1,\dots,n\} \to \N$ with $|\alpha|=3$,
let $v_\alpha := c_\alpha x^\alpha$ (multinomial coefficient). The
chain rule gives
\[
   \dot v_\alpha \;=\; \sum_s \alpha(s)\, c_\alpha\, \dot x_s\, x^{\alpha - e_s},
\]
where $e_s$ is the $s$-th standard basis vector. Substituting
$\dot x_s = \sum_{a,b} \widetilde q_{s,a,b} x_0^2 x_a x_b$ (post $r^2$-trick),
each positive entry $\widetilde q_{s,a,b} > 0$ contributes a
degree-$6$ positive monomial to $\dot v_\alpha$, with exponent
\begin{equation}
   \delta(k) \;:=\; \bigl(\alpha(k) - \iv{k=s}\bigr) + \mu_{0,a,b}(k).
   \tag{$\delta$}
\end{equation}
Provided $\alpha(s) > 0$, we have $\alpha(k) - \iv{k=s} \ge 0$ and the
expression is well-defined.

\subsection*{3c. The key combinatorial notion --- the split}
\label{sec:per-pair}

Given $\delta$ as above, a \emph{balanced split} is a pair of exponent
vectors $(\beta, \gamma)$ with
\[
   \beta + \gamma = \delta, \qquad |\beta| = |\gamma| = 3.
\]
A balanced split is \emph{non-autocatalytic for $\alpha$} if
\[
   \beta \ne \alpha \quad\text{and}\quad \gamma \ne \alpha.
\]
The interpretation is a cubic reaction $\beta \to \gamma$ whose reactant
and product triples are both different from $\alpha$ --- so $\alpha$
genuinely transitions and is not merely being catalyzed.

\medskip

\textbf{Why this definition?} We want a splitting that corresponds to
a reaction in which no $\alpha$-monomer is present on \emph{both}
sides. If $\beta = \alpha$, the reaction is $\alpha \to \gamma$ with
$\alpha$ present on the reactant side in full; but $\delta - \alpha =
\gamma$, meaning one copy of $\alpha$ just sits around unchanged during
the reaction. That is the classic self-catalysis pattern we are
excluding.

\section{Step 4: existence of a non-autocatalytic split, per pair}

This is where the work is. The theorem we prove says: whenever
$\widetilde q_{s,a,b} > 0$ with the slack structure in place,
a balanced non-autocatalytic split of $\delta$ exists.

\begin{theorem}[per-pair splitter]
\label{thm:per-pair}
Let $\widetilde q$ be slack-structured on species $\{0,1,\dots,n\}$, and
let $\alpha$ be a cubic exponent vector. Let $s, a, b$ be species with
$\alpha(s) > 0$, $a \ne 0$, $b \ne 0$ (so that $\widetilde q_{s,a,b}$ is
a positive \emph{non-slack-reactant} entry by (NSM)). Suppose
\begin{equation}
   s \ne 0 \;\Longrightarrow\; a \ne s \;\text{or}\; b \ne s. \tag{$\star$}
\end{equation}
Then $\delta$ (defined by ($\delta$)) admits a balanced non-autocatalytic
split.
\end{theorem}

We split the argument into three cases based on the identity of $s$ and
the location of $\alpha$'s mass.

\paragraph{Case A: $s = 0$ (slack is the source).}
Then the $-\iv{k=s}$ term in ($\delta$) subtracts a slack unit, and
$\mu_{0,a,b}$ adds two. Net effect: $\delta$ has at least one slack unit
and at least one unit each on $a$ and $b$ (which are non-slack by
hypothesis). So the support of $\delta$ contains $\{0, a, b\}$ with
multiplicities $\ge 1$ everywhere and the total is $|\delta| = 6$.

We split by:
\[
   \beta := e_0 + e_a + e_t, \qquad \gamma := e_0 + e_b + e_{t'},
\]
where $e_t, e_{t'}$ are chosen to make $\beta + \gamma = \delta$ and
to avoid equality with $\alpha$. Whenever $\delta - e_0 - e_a - e_b$
has support $\{0\}$ (i.e.\ $\delta$ concentrates on slack and $\{a,b\}$),
$\beta, \gamma$ trivially differ from $\alpha$ because $\alpha(0)$ would
have to absorb all three mass units of $\alpha$, but $\alpha$ also has
$\alpha(s) = \alpha(0) > 0$ and two more units distributed somewhere.
A case distinction on whether those remaining units fall on $a$, $b$,
or a third species completes the construction; two symmetric branches
(``distinct'' and ``doubled'') handle the sub-cases.

In all sub-cases $\beta \ne \alpha$ and $\gamma \ne \alpha$ because each
contains one slack unit, whereas $\alpha$ has at most $3$ slack units,
and the remaining structure differs.

\paragraph{Case B: $s \ne 0$ and $a \ne s$.}
This is the central case. The slack $0$ in $\delta$ enters with
multiplicity $2$ (from $\mu_{0,a,b}$). The source index $s$ enters with
multiplicity $\alpha(s) - 1 \ge 0$. The species $a$ enters with
multiplicity $\alpha(a) + 1 \ge 1$.

The \emph{foreign pair} is $(0, a)$: neither coordinate equals $s$. We
set
\[
   \beta := e_0 + e_a + \eta_\beta, \qquad \gamma := e_0 + e_\gamma + \eta_\gamma,
\]
where the concrete shapes of $\eta_\beta, \eta_\gamma$ are adjusted so
that $\beta + \gamma = \delta$ and the total on each side is $3$.

Why is $\beta \ne \alpha$? Because $\alpha(s) > 0$, if we had
$\beta = \alpha$, then the coordinate $s$ in $\beta$ would need to equal
$\alpha(s) \ge 1$. But $\beta$ only puts mass on $\{0, a, \text{one
more species}\}$, none of which is $s$ (since $a \ne s$, and any third
species we can always pick to differ from $s$). Contradiction.
Symmetrically for $\gamma$.

\medskip

The construction of the third species filling the remaining slot is a
combinatorial flow-matching step. The key observation is that we have
six units of mass in $\delta$ (two on slack, then the remainder), and
we need to split them into two halves of three, each avoiding $\alpha$.
This is a degree-3 bin-packing problem with finite feasibility; a direct
case analysis on the support sizes $|\supp(\alpha)| \in \{1, 2, 3\}$
shows a solution always exists. (This is the ``flow-matching lemma'',
which we sketch in \S\ref{sec:flow} below.)

\paragraph{Case C: $s \ne 0$ and $b \ne s$.}
Symmetric to Case B, with foreign pair $(0, b)$.

\medskip

The disjunction ($\star$) ensures that in the $s \ne 0$ branch at least
one of $a \ne s$, $b \ne s$ holds, so either Case B or Case C applies.

\begin{lemma}[($\star$) is automatic under slack structure]
\label{lem:star-auto}
If $\widetilde q$ satisfies (NSS), $\widetilde q_{s,a,b} > 0$, and
$s \ne 0$, then $a \ne s$ or $b \ne s$.
\end{lemma}
\begin{proof}
Suppose both $a = s$ and $b = s$. Then $\widetilde q_{s,a,b} =
\widetilde q_{s,s,s} \le 0$ by (NSS), contradicting positivity.
\end{proof}

Combining Theorem~\ref{thm:per-pair} with Lemma~\ref{lem:star-auto}:

\begin{corollary}[main technical result]
\label{cor:quad-split}
If $\widetilde q$ is a slack-structured QCF, then for every cubed
$\alpha$ and every positive entry $\widetilde q_{s,a,b} > 0$ with
$a, b \ne 0$ and $\alpha(s) > 0$, the monomial $\delta$ admits a
balanced non-autocatalytic split.
\end{corollary}

\section{Step 5: putting it together --- PP to NAP}

\begin{theorem}[PP~$\to$~NAP]
\label{thm:main}
Let $c_{r,i,j}$ be a PP on $n$ species. Fix $j_0 \in \{1,\dots,n\}$ and
$\lambda \in (0,1)$, and form the $\lambda$-lift $\widetilde q$ of the
folded tensor $q_{r,i,j} = c_{r,i,j} - \iv{i=r} - \iv{j=r}$ according to
$(\Lambda)$. Then for every cubed exponent vector $\alpha$ on $n+1$
species, every source $s$ with $\alpha(s) > 0$, and every pair $(a,b)$
with $a, b \ne 0$ and $\widetilde q_{s,a,b} > 0$: the associated
degree-$6$ monomial $\delta$ (after the $r^2$-trick) admits a balanced
non-autocatalytic split of $\alpha$.
\end{theorem}
\begin{proof}
Compose the three steps:
\begin{enumerate}
\item $c \ge 0$, $\sum_r c_{r,i,j} = 2$ $\Rightarrow$ $q$ satisfies
(NSS) (Lemma~\ref{lem:nss}).
\item $q$ is (NSS) and $\lambda \in (0,1)$ $\Rightarrow$ $\widetilde q$
is slack-structured on $n+1$ species (Lemmas~\ref{lem:lift-cons} and
\ref{lem:lift-slack}).
\item $\widetilde q$ is slack-structured $\Rightarrow$ for every
positive entry $\widetilde q_{s,a,b} > 0$ with $a,b \ne 0$, $\delta$
admits a balanced non-autocatalytic split for every $\alpha$ with
$\alpha(s) > 0$ (Corollary~\ref{cor:quad-split}).
\end{enumerate}
\end{proof}

\section{Aggregation and the reading as a CRN}

Theorem~\ref{thm:main} produces a split for each positive $(s,a,b)$
entry of $\widetilde q$ and each cubed $\alpha$ with $\alpha(s) > 0$.
The protocol-level picture follows by summing over these:

\begin{quote}
For every cubed $\alpha$ and every source $s$ with $\alpha(s) > 0$,
the $r^2$-lifted chain-rule contribution to $\dot v_\alpha$,
\[
   \alpha(s) \sum_{a,b} \widetilde q_{s,a,b}\, x^\delta,
\]
routes through a weighted sum of cubic reactions $\beta \to \gamma$,
none of which has $\alpha$ on either side.
\end{quote}

The routing uses a simple identity: for each positive entry, the rate
$\alpha(s)\,\widetilde q_{s,a,b}$ is assigned to the split reaction
$\beta \to \gamma$ found by Theorem~\ref{thm:per-pair}. The sum over
$(a,b)$ with $\widetilde q_{s,a,b} > 0$ is the total positive
contribution.

\emph{What about negative entries?} A negative entry
$\widetilde q_{s,a,b} < 0$ corresponds, in the PP picture, to the
folding terms $-\iv{i=r} - \iv{j=r}$: these are \emph{degradation}
edges attached to the reaction ``$(i,j) \to$ (something)''. They are not
independent free variables but consumption terms coupled to the same
reaction via (P4). In the CRN translation, they are realized as the
loss of one copy of $i$ (or $j$) upon firing, which is structurally
non-autocatalytic: destroying one copy of a species is obviously not
self-replication.

\section{Why the forbidden-edge concern does not bite}
\label{sec:forbidden}

A natural worry is the following. Suppose a monomial $-k\, u\, v$ shows
up in $\dot v$ (negative coefficient, large $k$). The same monomial
$u\, v$ can also appear with a positive coefficient $K$ in some
$\dot w$. If $k$ is large enough, wouldn't removing the forbidden edge
$(-k)$ destroy the feasibility of a global flow condition and invalidate
the split?

The answer is that no such global flow condition appears in our
argument. Three layers of structural isolation rule it out:

\paragraph{First layer: folding-derived negatives are reaction bookkeeping,
not free edges.} The only negatives appearing in our QCF are of the
form $-\iv{i=r} - \iv{j=r}$, which record ``when $i$ and $j$ react, one
copy of $i$ is consumed, and so is one copy of $j$''. This is the same
reaction already accounted for on the positive side, not a free edge
to match independently.

\paragraph{Second layer: the coefficients of a single monomial across
different $\dot x_r$ are tied by (P1).} In the original PP, every
monomial $x_i x_j$ corresponds to one reaction $i + j \to \text{products}$,
with the distribution over product species governed by
$c_{r,i,j}/2$ and $\sum_r c_{r,i,j} = 2$. The ``$-k$'' and ``$+K$'' the
worry imagines are not independent parameters: they are both extracted
from the \emph{same} probability vector $c_{\cdot, u, v}$. You cannot
have ``large $k$'' without forcing the corresponding $c$'s to be
small; the coupling is through (P1). Thus a large forbidden edge
\emph{cannot} be created in isolation of its positive counterparts.

\paragraph{Third layer: the splitting is per pair, not global.}
The per-pair theorem (Theorem~\ref{thm:per-pair}) takes one $(s,a,b)$
at a time and produces a split using only \emph{that entry}'s
positivity and the slack structure. There is no cross-$(s,a,b)$ flow
feasibility that has to be solved. The $r^2$-trick contributes one
slack factor of multiplicity $2$ per entry; the original PP
contributes at least one non-source reactant (from (NSS) via
Lemma~\ref{lem:star-auto}); these are enough to guarantee the
foreign-pair condition ($\star$) independently for each entry.
Aggregation over all entries is then a linear sum --- the rates for
each positive entry add up --- with no constraint from the negatives.

\medskip

In short: the proof never asks ``is there a feasible flow routing all
positive and negative edges?'' It asks ``given each positive edge, can
I reshape its cubic footprint into a non-autocatalytic transition?''
The answer is yes, one edge at a time; the negatives are already
bookkeeping on the same reaction.

\section{The flow-matching lemma}
\label{sec:flow}

For completeness we sketch the degree-$3$ existence claim used inside
Case B and Case C of Theorem~\ref{thm:per-pair}.

\begin{lemma}[split existence given foreign pair]
\label{lem:flow}
Let $\delta : \{0,1,\dots,n\} \to \N$ with $|\delta|=6$. Let $\alpha$ be
a cubic exponent vector. Suppose there exist two distinct indices
$i_1, i_2 \in \supp(\delta)$ with
\[
   i_1 \ne s, \qquad i_2 \ne s, \qquad \delta(i_1) \ge 1, \qquad \delta(i_2) \ge 1,
\]
where $s$ is a fixed ``source'' species with $\alpha(s) > 0$. Then
there is a balanced split $(\beta, \gamma)$ of $\delta$ with
$\beta \ne \alpha$ and $\gamma \ne \alpha$.
\end{lemma}

\begin{proof}[Proof sketch]
Start by placing $e_{i_1} + e_{i_2}$ into $\beta$; this already ensures
$\beta \ne \alpha$ because $\beta$ now contains at least one of
$i_1, i_2 \ne s$, whereas the triple $\alpha$ (with $\alpha(s) \ge 1$)
has at least one unit at $s$ that $\beta$ lacks.

The remaining mass is $\delta - e_{i_1} - e_{i_2}$, of total $4$. We need
to add one more unit to $\beta$ (to reach size $3$) and the rest to
$\gamma$ (of size $3$). Choose the single unit for $\beta$ so that
$\gamma = \delta - \beta$ also differs from $\alpha$. Two sub-cases:
\begin{itemize}
\item If $\delta - e_{i_1} - e_{i_2}$ places mass on a species $t \ne s$,
take that $t$ to complete $\beta$; then $\gamma$'s remaining mass on
$s$ is $\delta(s) = \alpha(s)$ minus whatever adjustments --- and the
supports disagree with $\alpha$ because $\gamma$'s third unit is at some
$t' \ne s$ (by residual mass count, $\gamma$ always has at least one
non-$s$ unit).
\item If the remaining mass lives entirely on $s$, then we use the fact
that $|\supp(\alpha)| \le 3$ but $\gamma$'s support of size $3$ then
concentrates at $s$ with multiplicity $3$, which forces
$\alpha = 3 e_s$ for equality. In that case we re-route one unit from
$\beta$ to $\gamma$ and rebalance (swapping one unit between $\beta$
and $\gamma$) --- again possible because at least two distinct
off-$s$ indices are available.
\end{itemize}
In each sub-case explicit formulas for $(\beta, \gamma)$ can be
written down; the verification that $\beta, \gamma \ne \alpha$ is a
one-line count of multiplicities.
\end{proof}

\begin{remark}
The $r^2$-trick guarantees $\delta(0) \ge 2$, so the slack $0$ can
always serve as one of the indices $i_1, i_2$. The other index comes
from the PP reactant distinction: if $s \ne 0$, the non-autocatalytic
hypothesis ($\star$) yields either $a \ne s$ or $b \ne s$, providing
the second foreign index.
\end{remark}

\section{Summary}

The argument is modular. Three independent steps compose into the
main theorem:

\begin{center}
\begin{tabular}{r p{0.7\textwidth}}
\textbf{Step 1.} & Fold the PP into a signed conservative field $q$
satisfying (NSS). \\[2pt]
\textbf{Step 2.} & $\lambda$-lift into $\widetilde q$ on one extra slack
species, preserving (NSS) and gaining (NSM). \\[2pt]
\textbf{Step 3.} & $r^2$-trick and cubing: every positive entry of
$\widetilde q$ together with any cubic $\alpha$ with non-zero source
mass produces a degree-$6$ footprint $\delta$ admitting a balanced
non-autocatalytic split.
\end{tabular}
\end{center}

No step requires a global flow-feasibility check; each entry is
handled in isolation. The negatives introduced by folding are reaction
bookkeeping, not free edges, and they are tied by (P1) to the positive
entries they accompany. This is why the ``forbidden edge with large
weight'' scenario does not break the construction.

\end{document}
